\begin{abstract}
Depth-pruned language models are often continued-pretrained until held-out
perplexity improves, implicitly treating likelihood recovery as evidence that
target capabilities have recovered. We test that proxy on literal observed
paths from OLMo-2-7B, using in-domain perplexity, two MMLU scoring interfaces,
and three zero-shot closed-book QA evaluations. In the single principal
keep14+fresh2 run, both perplexity and answer-letter MMLU improve through the
available 200k-step path, but the endpoint remains far below the intact model on
MMLU and all three QA tasks. The intact full32 branch is observed only through
25k; it stays closer to the base than keep14 at that short horizon, but does not
control the 200k endpoint. Complete-option
MMLU narrows the gap; however, a fully random 16-layer operating point reaches a
similar content-score floor while remaining at answer-letter chance. A
ShortGPT-16 operating point is substantially stronger, but its construction
couples block selection, inherited-block count, final-block retention, and the
absence of fresh tails. Thus, PPL improvement alone does not imply target
recovery along these observed paths. We recommend reporting in-domain
likelihood, target evaluation, interface, exact
construction, training budget, and run-level uncertainty separately. We do not
claim knowledge deletion or localization, causal factor attribution, universal
recovery dynamics, or failure beyond the measured budgets.
\end{abstract}
